% ============================================================================================
% BAB III ANALISIS MASALAH
% Pembagian subbab tidak rigid dan dapat bervariasi. Bab ini minimal berisi analisis kebutuhan
% fungsional dan nonfungsional, analisis berbagai alternatif solusi yang dapat ditawarkan, dan
% metode pemilihan solusi yang diusulkan.
% ============================================================================================

\chapter{ANALISIS MASALAH}
\label{chap:analisis-masalah}
\section{Analisis Kondisi Saat Ini}
\label{sec:kondisi_saat_ini}

% Analisis terhadap lanskap teknologi pembelajaran adaptif saat ini mengungkapkan kesenjangan fundamental antara kemampuan sistem untuk mempersonalisasi materi dan ketidakmampuannya untuk mengelola dinamika sosial pembelajar. Kondisi ini dapat dipetakan ke dalam tiga masalah arsitektural utama.

% \subsection{Dominasi Algoritma Berbasis Isolasi}
Mayoritas sistem asesmen adaptif yang ada saat ini dibangun di atas asumsi bahwa pembelajaran adalah proses soliter. Algoritma pelacakan pengetahuan (\textit{knowledge tracing}) yang dominan bekerja dengan memproses data interaksi biner (benar/salah) antara siswa dan mesin secara eksklusif sesuai dengan yang telah dibahas pada subbab \ref{sec:teknik_algoritma_aa}.

Implikasi dari desain ini adalah terjadinya \textit{algorithmic blindness} terhadap konteks sosial. Sistem mampu memodelkan kemahiran individu dengan presisi tinggi, namun tidak dapat melihat fakta bahwa siswa mungkin terjebak dalam stagnasi pembelajaran \textit{learning plateau} atau \textit{filter bubble} akademik. Pelajar terus menerus hanya mengerjakan tipe soal yang dikuasai dan tidak pernah terekspos pada perspektif orang lain \autocite{Halkiopoulos2024}. Akibatnya, sistem berisiko untuk justru menghambat perkembangan kognitif karena tidak mampu mengeluarkan individu dari \textit{learning plateau} hanya dengan interaksi satu arah dari sistem yang terus menerus mengakomodasi tipe soal yang dikuasai oleh siswa.

% \subsection{Fragmentasi Antara Asesmen Kognitif dan Aktivitas Sosial}
Dalam kondisi saat ini, aktivitas sosial sering kali hanya ditempatkan sebagai fitur tambahan (\textit{add-on}) yang berjalan paralel namun terpisah dari mesin asesmen utama. Data berharga dari interaksi sosial seperti kemampuan siswa dalam mengevaluasi pekerjaan orang lain atau memberikan argumen seringkali terbuang dan tidak dikonversi menjadi sinyal kompetensi yang dapat mengubah rekomendasi materi \autocite{Ihichr2024}.

% \subsection{Keterbatasan Mitigasi Berbasis Kendali Pengguna}
Upaya yang ada untuk mengatasi masalah isolasi pembelajaran sejauh ini cenderung mengandalkan inisiatif pengguna, seperti pendekatan \textit{User Controllable Recommender Systems} (UCRS) seperti yang dibahas pada subbab \ref{sec:upaya_penyelesaian} atau visualisasi sosial (OSSM) yang dibahas pada subbab \ref{sec:upaya_integrasi}. Namun, pendekatan ini memiliki kelemahan fundamental dalam konteks pendidikan, pendekatan ini berasumsi bahwa pengguna menyadari bahwa mereka sedang terisolasi sehingga secara mandiri akan melakukan mitigasi melalui fitur yang disediakan.

Dalam kenyataannya, fenomena \textit{overpersonalization} yang terealisasi dengan gejala \textit{filter bubble} atau \textit{echo-chamber} seringkali bersifat tidak terlihat. Sistem secara otomatis menyaring konten yang tidak sesuai dengan profil historis sehingga siswa tidak pernah terekspos pada opsi alternatif dan menganggap bahwa materi yang mereka terima adalah satu-satunya yang tersedia. Oleh karena itu, menyerahkan kendali mitigasi sepenuhnya kepada pengguna berisiko tidak efektif karena pengguna tidak dapat memilih apa yang tidak mereka lihat. Sistem saat ini kekurangan mekanisme intervensi otomatis yang mampu memaksa terjadinya diversifikasi interaksi ketika pola stagnasi terdeteksi oleh algoritma.
% \subsection{Implikasi Terhadap Arah Pengembangan Sistem}
Dari analisis kondisi di atas, dapat disimpulkan bahwa menambahkan fitur kolaborasi konvensional tidak akan menyelesaikan masalah \textit{overpersonalization} selama mesin asesmennya masih bekerja secara terisolasi.

Kebutuhan mendesak saat ini adalah melakukan rekayasa ulang (\textit{re-engineering}) pada arsitektur asesmen adaptif. Diperlukan sebuah sistem yang mampu mengintegrasikan aktivitas sosial bukan sekadar sebagai fitur tambahan, melainkan sebagai instrumen intervensi proaktif yang menyatu dengan pelacakan kompetensi individu. Sistem harus memiliki kapabilitas untuk: (1) mendeteksi tanda-tanda stagnasi kognitif (\textit{learning plateau}) akibat personalisasi berlebih secara otomatis dan (2) merekomendasikan interaksi sosial yang terarah secara heterogen untuk menyisipkan keberagaman perspektif untuk mendorong pelajar keluar dari stagnasi kognitif. Dengan arsitektur yang interaktif dan dinamis ini, sistem dapat memitigasi risiko menghambat di \textit{learning plateau} secara sistematis tanpa mengabaikan integritas penilaian kognitif individu.

\section{Analisis Kebutuhan}
\label{sec:analisis_kebutuhan}

Berdasarkan analisis kondisi saat ini pada subbab \ref{sec:kondisi_saat_ini}, sistem yang dibutuhkan harus mampu menjembatani keterpisahan antara personalisasi individu dengan lingkungan sosialnya untuk memitigasi risiko \textit{overpersonalization}. Kebutuhan ini dirumuskan tanpa membatasi teknik implementasi spesifik, guna memungkinkan eksplorasi berbagai alternatif solusi (seperti \textit{peer assessment}, \textit{mentoring}, atau rekomendasi jejaring) pada tahap pemilihan solusi.

\subsection{Kebutuhan Fungsional}
\label{ssec:kebutuhan_fungsional}
Daftar kebutuhan fungsional sistem dari sisi pengguna (siswa) dan logika sistem disajikan pada Tabel \ref{tbl:kebutuhan-fungsional}.

\begin{table}[H]
    \centering\caption{Kebutuhan Fungsional (FR) Sistem}\label{tbl:kebutuhan-fungsional}
% Menambah sedikit jarak antar baris agar tidak terlalu padat
\renewcommand{\arraystretch}{1.3}
\begin{tabular}{|l|p{0.85\textwidth}|}
\hline
\textbf{Kode} & \textbf{Deskripsi Kebutuhan Fungsional} \\
\hline

% Bagian 1
\multicolumn{2}{|l|}{\textbf{Interaksi Pengguna (pemelajar)}} \\
\hline
FR-01 & Sistem harus menyediakan mekanisme otentikasi bagi pemelajar untuk mengakses lingkungan pembelajaran personal mereka. \\
\hline
FR-02 & Sistem harus menyajikan konten pembelajaran, spesifiknya konten asesmen, yang telah disesuaikan tingkat kesulitannya dengan kemahiran pemelajar saat ini. \\
\hline
FR-03 & Sistem harus memfasilitasi \textbf{interaksi pembelajaran antar-pengguna} yang memungkinkan terjadinya pertukaran informasi atau kolaborasi akademik di dalam platform. \\
\hline
FR-04 & Sistem harus menyediakan informasi mengenai konteks sosial pembelajaran (misalnya, aktivitas komunitas atau capaian kolektif) untuk membangun kesadaran sosial (\textit{social awareness}). \\
\hline

% Bagian 2
\multicolumn{2}{|l|}{\textbf{Logika Adaptasi \& Pemrosesan Data (Sistem)}} \\
\hline
FR-05 & Sistem harus memiliki algoritma untuk mengestimasi tingkat kemahiran pengetahuan individu pemelajar secara dinamis berdasarkan riwayat aktivitasnya. \\
\hline
% FR-06 & Sistem harus mampu \textbf{mengkuantifikasi data interaksi sosial} (dari FR-03) menjadi parameter numerik yang dapat diproses lebih lanjut (misalnya, skor keterlibatan atau skor kinerja kolaboratif). \\
FR-06 & Sistem harus mampu mendeteksi kondisi stagnasi pembelajaran (\textit{learning plateau}) pemelajar secara otomatis berdasarkan suatu metrik atau proksi. \\
\hline
FR-07 & Sistem harus mampu merekomendasikan \textbf{intervensi sosial spesifik} (misalnya, merekomendasikan orang untuk berinteraksi atau topik untuk didiskusikan) agar bisa mendorong pemelajar keluar dari stagnasi pembelajaran. \\
\hline
FR-08 & Sistem harus menggunakan metrik kuantitatif untuk menentukan materi atau aktivitas berikutnya yang paling optimal bagi pemelajar. \\
\hline
FR-09 & Sistem harus memiliki mekanisme untuk menggabungkan alur pembelajaran yang adaptif secara individual (FR-05 dan FR-08) dengan aktivitas sosial (FR-07) untuk memberikan diversifikasi konten pembelajaran. \\
\hline
\end{tabular}
\end{table}

\subsection{Kebutuhan Nonfungsional}
\label{ssec:kebutuhan_nonfungsional}
Kebutuhan nonfungsional disesuaikan untuk mendukung interaksi sosial yang responsif tanpa menentukan metode spesifik.
Tabel \ref{tbl:kebutuhan-nonfungsional} merinci kebutuhan nonfungsional sistem yang diusulkan.

\begin{table}[H]
\centering
\caption{Kebutuhan Nonfungsional (NFR) Sistem}\label{tbl:kebutuhan-nonfungsional}
\renewcommand{\arraystretch}{1.3}
\begin{tabular}{|l|p{0.8\textwidth}|}
\hline
\textbf{Kode} & \textbf{Deskripsi Kebutuhan Nonfungsional} \\
\hline
NFR-01 & \textbf{(Usability)} Antarmuka interaksi sosial harus terintegrasi secara intuitif dalam platform yang sama dengan platform pembelajaran, meminimalkan beban kognitif siswa saat berpindah antara belajar mandiri dan berkolaborasi. \\
\hline
NFR-02 & \textbf{(Interpretability)} Sistem harus mampu memberikan indikasi alasan rekomendasi (baik materi maupun interaksi sosial). \\
\hline
NFR-03 & \textbf{(Performance)} Pembaruan model pengguna (baik skor individu maupun sosial) harus terjadi dalam waktu nyata (\textit{near real-time}) setelah interaksi selesai dilakukan. \\
\hline
NFR-04 & \textbf{(Scalability)} Sistem harus mampu menangani konkurensi interaksi sosial dari satu kelas (±50 pengguna) secara simultan. \\
\hline
NFR-05 & \textbf{(Privacy)} Mekanisme interaksi sosial harus dirancang sedemikian rupa sehingga data privasi spesifik (seperti nilai asesmen mentah) tidak terekspos ke pengguna lain. \\
\hline
NFR-06 & \textbf{(Accessibility)} Sistem dapat diakses melalui peramban web standar. \\
\hline
\end{tabular}
\end{table}

\subsection{Pemetaan Kebutuhan}
\label{ssec:pemetaan_kebutuhan}
Tabel \ref{tbl:pemetaan-kebutuhan} memetakan tujuan penanganan masalah ke fungsi spesifik yang telah digeneralisasi.

\input{tables/tabel_3_pemetaan.tex}

\section{Analisis Pemilihan Solusi}
\label{sec:analisis_pemilihan_solusi}
Untuk menjawab kebutuhan sistem dalam memitigasi \textit{overpersonalization}, diperlukan pemilihan komponen teknologi yang tidak hanya akurat, tetapi juga fungsional untuk diintegrasikan. Pemilihan solusi dilakukan pada tiga dimensi utama: (1) teknik asesmen kemahiran individu, (2) teknik interaksi dan asesmen sosial, dan (3) teknik spesifik pencegahan \textit{echo chamber}.

Penentuan solusi teknis untuk setiap dimensi dilakukan secara objektif menggunakan metode \textit{Weighted Decision Matrix} (Matriks Keputusan Terbobot). Kriteria evaluasi diturunkan dari Batasan Masalah (Bab \ref{chap:pendahuluan}) dan Tinjauan Literatur (Bab \ref{chap:studi-literatur}).

\subsection{Penentuan Teknik Asesmen Individu}
\label{ssec:keputusan_asesmen_individu}
Komponen ini berfungsi memodelkan kompetensi siswa. Dalam konteks pencegahan \textit{overpersonalization}, kriteria utamanya adalah kemampuan model untuk melakukan estimasi tingkat kemahiran secara dinamis (\textit{near real-time}) setelah setiap aktivitas, serta menghasilkan trajektori kompetensi kontinu yang dapat dianalisis secara matematis untuk mendeteksi tanda-tanda stagnasi kognitif (\textit{learning plateau}).

Tinjauan literatur pada Subbab \ref{sec:teknik_algoritma_aa} mengidentifikasi lima alternatif utama untuk algoritma pelacakan pengetahuan, yaitu IRT/Bayesian IRT, FA, BKT, DKT, dan \textit{Elo rating}:

\begin{enumerate}
    \item \textbf{\textit{Item Response Theory} (IRT) dan \textit{Bayesian IRT} (IRT*):} Model IRT cenderung statis karena memerlukan sampel data yang besar untuk mengkalibrasi parameter butir soal sebelum dapat digunakan. Hal ini membuat IRT tidak praktis untuk melakukan estimasi ulang terhadap parameter kemahiran ($\theta$) secara dinamis setiap kali siswa menyelesaikan satu butir soal, sehingga menyulitkan deteksi stagnasi kognitif secara cepat.

    \item \textbf{\textit{Factor Analysis} (FA):} FA dirancang untuk validasi pasca-uji (\textit{post-hoc}) guna mengidentifikasi faktor laten secara kolektif, bukan pelacakan kemahiran dinamis per individu. Karena pembaruan parameter tidak terjadi secara langsung saat interaksi siswa berlangsung, FA tidak dapat digunakan sebagai basis data trajektori kontinu untuk mendeteksi stagnasi kognitif.

    \item \textbf{\textit{Bayesian Knowledge Tracing} (BKT):} BKT melacak kemahiran dalam bentuk status biner (belum belajar vs. sudah belajar) pada konsep keterampilan yang sangat spesifik. Karakteristik biner ini menyulitkan sistem untuk mendeteksi \textit{learning plateau} pada rentang kinerja yang kontinu, karena siswa mungkin sudah terjebak dalam stagnasi kognitif sebelum terjadi perubahan status penguasaan konsep di dalam model.

    \item \textbf{\textit{Deep Knowledge Tracing} (DKT):} Meskipun akurasi pelacakannya sangat tinggi, sifat DKT yang berupa \textit{black-box} membuat parameter internalnya sulit diinterpretasikan. Ketiadaan interpretabilitas ini menyulitkan sistem untuk merumuskan ambang batas variansi perubahan kompetensi yang logis dan terukur sebagai proksi deteksi stagnasi kognitif.

    \item \textbf{Sistem \textit{Elo rating}:} Sistem \textit{Elo rating} mampu memperbarui peringkat kemahiran secara instan setelah setiap jawaban dikirimkan tanpa memerlukan prakalibrasi soal. Trajektori skor kemahiran yang kontinu dan terbarui secara dinamis ini sangat ideal untuk menghitung nilai variansi perubahan skor ($\text{Var}(\Delta\theta)$) guna mendeteksi \textit{learning plateau} secara otomatis.
\end{enumerate}

Matriks keputusan untuk pemilihan teknik asesmen individu disajikan pada Tabel \ref{tbl:solusi-r1}.

\begin{table}[H]
    \centering
    \caption{Matriks Keputusan Teknik Asesmen Individu}
    \label{tbl:solusi-r1}
    \renewcommand{\arraystretch}{1.3} % Memberi sedikit ruang vertikal
    \begin{tabular}{|p{4.0cm}|c|c|c|c|c|}
    \hline
    \textbf{Kriteria} & \textbf{IRT} & \textbf{FA} & \textbf{BKT} & \textbf{DKT} & \textbf{Elo} \\
    \hline
    Efisiensi Komputasi & 2 & 1 & 3 & 1 & 5 \\
    \hline
    Kebutuhan Data Awal (\textit{Cold Start}) & 2 & 2 & 3 & 2 & 5 \\
    \hline
    Akurasi Pelacakan & 4 & 3 & 3 & 5 & 4 \\
    \hline
    Interpretabilitas & 4 & 3 & 4 & 1 & 5 \\
    \hline
    \textbf{Rata-rata Skor} & \textbf{3.00} & \textbf{2.25} & \textbf{3.25} & \textbf{2.25} & \textbf{4.75} \\
    \hline
    \end{tabular}
\end{table}

\textbf{Rasionalisasi Objektif dan Analisis \textit{Trade-off}:}
\begin{enumerate}
    \item \textbf{Efisiensi \& Kebutuhan Data (Bobot Tinggi - 60\%):} Kriteria ini diberi bobot tertinggi mengingat batasan masalah pengembangan purwarupa yang mungkin tidak memiliki dataset historis besar. \textbf{\textit{Elo rating}} unggul mutlak (skor 5) karena algoritma ini dirancang untuk bekerja tanpa prakalibrasi yang rumit dan efisien secara komputasi. Sebaliknya, \textbf{IRT} dan \textbf{FA} membutuhkan sampel besar untuk kalibrasi parameter butir soal, dan \textbf{DKT} membutuhkan data pelatihan masif untuk melatih jaringan saraf.
    
    \item \textbf{Akurasi vs. Kelayakan (\textit{Trade-off}):} \textbf{DKT} memiliki skor akurasi tertinggi (5) karena kemampuannya menangkap pola kompleks. Namun, terjadi \textit{trade-off} signifikan pada aspek \textbf{Interpretabilitas} (skor 1) karena sifat \textit{black-box}-nya, yang menyulitkan sistem untuk menjelaskan alasan rekomendasi kepada siswa. \textbf{\textit{Elo rating}} dipilih bukan karena merupakan teknik paling akurat, tetapi karena \textit{Elo rating} adalah solusi paling layak (\textit{feasible}) untuk diimplementasikan dalam skenario data terbatas dengan akurasi yang masih dapat diterima.
\end{enumerate}

\subsection{Penentuan Teknik Interaksi dan Asesmen Sosial}
\label{ssec:keputusan_asesmen_sosial}
Komponen ini bertujuan memfasilitasi interaksi yang bermakna. Untuk memecahkan \textit{filter bubble}, aktivitas sosial harus memaksa siswa terpapar pada perspektif atau materi di luar preferensi historis mereka.

Evaluasi dilakukan terhadap lima alternatif bentuk aktivitas sosial, yaitu PBL, Peer Mentoring, Group Formation, Open Discussion, dan Peer Assessment, sebagaimana dibahas pada Subbab \ref{sec:collaborative_learning}:

\begin{enumerate}
    \item \textbf{Pembelajaran Berbasis Proyek (\textit{Project-Based Learning} / PBL):} Sangat baik untuk kolaborasi mendalam, namun memiliki kendala besar dalam penilaian otomatis. Sulit bagi sistem untuk memisahkan kontribusi individu dari hasil kelompok (\textit{free-rider problem}) secara akurat tanpa intervensi manual dosen sehingga sulit dijadikan input data \textit{real-time} untuk algoritma adaptif.

    \item \textbf{Pendampingan Sejawat (\textit{Peer Mentoring}):} Hubungan ini seringkali bersifat hierarkis dan searah (senior mengajar junior), bukan kolaborasi setara. Hal ini kurang efektif untuk tujuan memecahkan \textit{echo chamber} yang membutuhkan paparan terhadap berbagai perspektif sejawat yang beragam, bukan sekadar transfer pengetahuan vertikal.

    \item \textbf{Formasi Kelompok Berbasis Klastering (\textit{Group Formation}):} Efektif untuk "memaksa" keberagaman di awal (\textit{cold start}). Namun, karena sifatnya statis, teknik ini tidak menghasilkan data kinerja kontinu yang bisa digunakan sistem untuk memantau apakah siswa kembali mengisolasi diri setelah kelompok terbentuk.

    \item \textbf{Sesi Diskusi Terbuka (\textit{Open Discussion}):} Dalam forum bebas, siswa cenderung mencari teman yang sepemikiran (\textit{homophily}), yang justru memperkuat \textit{echo chamber} alih-alih memecahkannya. Selain itu, sulit menilai kualitas interaksi secara otomatis tanpa NLP yang kompleks.

    \item \textbf{Penilaian Sejawat (\textit{Peer Assessment}):} Teknik ini paling potensial karena sistem memegang kendali penuh atas "siapa menilai siapa". Sistem dapat menugaskan siswa untuk menilai rekan yang memiliki sudut pandang bertentangan atau tingkat kemampuan berbeda (di luar \textit{bubble}-nya). Selain itu, kualitas umpan balik dapat diukur sebagai metrik keterbukaan siswa terhadap kolaborasi.
\end{enumerate}

Matriks keputusan untuk pemilihan teknik interaksi sosial disajikan pada Tabel \ref{tbl:solusi-r2}.

\begin{table}[H]
    \centering
    \caption{Matriks Keputusan Teknik Interaksi Sosial}
    \label{tbl:solusi-r2}
    \renewcommand{\arraystretch}{1.3} % Menambah jarak antar baris
    \begin{tabular}{|p{4.0cm}|c|c|c|c|c|}
    \hline
    \textbf{Kriteria} & \textbf{PBL} & \textbf{Mentor} & \textbf{Group} & \textbf{Diskusi} & \textbf{Peer Assmt.} \\
    \hline
    Keterukuran Data (Measurability) & 2 & 3 & 1 & 1 & 5 \\
    \hline
    Kesetaraan Interaksi & 5 & 2 & 5 & 5 & 5 \\
    \hline
    Potensi Eksposur Keberagaman & 4 & 2 & 3 & 3 & 5 \\
    \hline
    Kompleksitas Implementasi & 2 & 3 & 4 & 3 & 4 \\
    \hline
    \textbf{Rata-rata Skor} & \textbf{3.25} & \textbf{2.50} & \textbf{3.25} & \textbf{3.00} & \textbf{4.75} \\
    \hline
    \end{tabular}
\end{table}

\textbf{Rasionalisasi Objektif dan Analisis \textit{Trade-off}:}
\begin{enumerate}
    \item \textbf{Keterukuran Data (35\%):} Sistem membutuhkan input kuantitatif untuk algoritma (FR-06). \textit{Peer Assessment} mendapatkan skor tertinggi (5) karena menghasilkan artefak terstruktur (skor rubrik dan teks umpan balik) yang mudah dikuantifikasi. Sebaliknya, \textbf{Diskusi Terbuka} dan \textbf{PBL} (Project-Based Learning) sangat sulit dinilai kontribusi individunya secara otomatis tanpa intervensi manual dosen.
    
    \item \textbf{Kesetaraan \& Eksposur (\textit{Trade-off}):} \textbf{Mentoring} mendapat skor rendah pada kesetaraan (2) karena sifatnya hierarkis (senior-junior), yang kurang cocok untuk memecah \textit{echo chamber} dibandingkan interaksi setara. \textbf{PBL} sebenarnya sangat baik dalam kesetaraan dan kolaborasi mendalam (skor 5), namun dipilihnya \textbf{Peer Assessment} merupakan \textit{trade-off} logis karena PBL memiliki kompleksitas implementasi yang tinggi dan risiko \textit{free-rider} yang sulit dideteksi sistem. Peer Assessment dipilih karena memungkinkan sistem mengatur siapa menilai siapa sehingga bisa memaksa eksposur keberagaman yang dibutuhkan.
\end{enumerate}

\subsection{Penentuan Teknik Mitigasi \textit{Overpersonalization}}
\label{ssec:keputusan_mitigasi_bubble}
Keputusan ini bertujuan memilih mekanisme teknis untuk memecahkan \textit{echo chamber} yang dapat beroperasi secara otomatis dan layak diimplementasikan secepat mungkin tanpa data pelatihan (\textit{cold-start}).

Sistem mengevaluasi empat alternatif teknik yang dirancang untuk mencegah polarisasi opini dan isolasi informasi, yaitu UCRS, GRS, FRediECH, dan \textit{Constraint-based Re-ranking}, sebagaimana dibahas pada Subbab \ref{sec:upaya_penyelesaian}:

\begin{enumerate}
    \item \textbf{\textit{User Controllable Recommender Systems} (UCRS):} Pendekatan ini sangat bergantung pada kesadaran pengguna. Dalam konteks pendidikan, siswa seringkali tidak menyadari bahwa mereka berada dalam \textit{bubble} kompetensi ("\textit{unconscious incompetence}"). Mengandalkan siswa untuk secara manual meminta materi yang mereka tidak sukai berisiko tidak efektif dibandingkan intervensi otomatis.

    \item \textbf{\textit{Group Recommender Systems} (GRS):} GRS cenderung memprioritaskan konsensus kelompok atau rata-rata preferensi. Dalam pembelajaran adaptif, hal ini berisiko mengorbankan personalisasi spesifik yang dibutuhkan siswa \textit{outlier} (sangat mahir atau sangat tertinggal) karena rekomendasi akan tertarik ke arah "pusat" kelompok, yang bisa menyebabkan \textit{under-personalization}.

    \item \textbf{Pemodelan Komunitas Implisit (\textit{FRediECH}):} Teknik ini sangat bagus untuk memecah polarisasi opini. Namun, kelemahan utamanya adalah kebutuhan akan data pelatihan yang besar untuk membangun graf interaksi awal (\textit{cold-start problem}) sehingga sulit diterapkan pada fase awal implementasi sistem tanpa dataset historis.

    \item \textbf{\textit{Constraint-based Re-ranking}:} Pendekatan ini bersifat agnostik terhadap model dasar dan bekerja berdasarkan aturan logis ($\alpha$ dan $\beta$) pada daftar rekomendasi yang ada. Keunggulan utamanya adalah efisiensi komputasi dan kemampuan untuk langsung berfungsi tanpa memerlukan pelatihan model data yang ekstensif (\textit{training-free}), menjadikannya solusi yang sangat potensial untuk kondisi data terbatas.
\end{enumerate}

Matriks keputusan untuk pemilihan teknik mitigasi disajikan pada Tabel \ref{tbl:solusi-r3}.

\begin{table}[H]
    \centering
    \caption{Matriks Keputusan Teknik Mitigasi Isolasi}
    \label{tbl:solusi-r3}
    \renewcommand{\arraystretch}{1.3} % Menambah spasi vertikal agar rapi
    \begin{tabular}{|p{4.5cm}|c|c|c|c|}
    \hline
    \textbf{Kriteria} & \textbf{(UCRS)} & \textbf{(GRS)} & \textbf{(FRediECH)} & \textbf{(Re-rank)} \\
    \hline
    Automasi Intervensi & 1 & 4 & 5 & 5 \\
    \hline
    Kelayakan \textit{Cold-Start} & 5 & 4 & 1 & 5 \\
    \hline
    Preservasi Personalisasi & 5 & 2 & 4 & 4 \\
    \hline
    \textbf{Rata-rata Skor} & \textbf{3.67} & \textbf{3.33} & \textbf{3.33} & \textbf{4.67} \\
    \hline
    \end{tabular}
\end{table}

\textbf{Rasionalisasi Objektif dan Analisis \textit{Trade-off}:}
\begin{enumerate}
    \item \textbf{Automasi Intervensi (35\%):} \textbf{UCRS} mendapat skor terendah (1) karena bergantung pada inisiatif pengguna yang seringkali bias atau pasif. \textbf{FRediECH} dan \textbf{Re-ranking} unggul (5) karena sistem secara proaktif melakukan intervensi untuk memecah \textit{bubble} tanpa menunggu tindakan pengguna.
    
    \item \textbf{Kelayakan \textit{Cold-Start} (35\%):} Kriteria ini krusial untuk purwarupa. \textbf{FRediECH} mendapat penalti besar (1) karena berbasis \textit{Deep Learning} yang membutuhkan data historis masif untuk pelatihan awal. Sebaliknya, \textbf{Re-ranking} mendapat skor sempurna (5) karena algoritma ini bekerja secara \textit{greedy} pada daftar yang ada berdasarkan aturan logis sehingga bisa langsung berjalan sejak hari pertama implementasi.
    
    \item \textbf{Preservasi Personalisasi (30\%):} \textbf{GRS} cenderung menarik rekomendasi ke rata-rata kelompok, mengorbankan preferensi unik individu (skor 2). \textbf{Re-ranking} dinilai baik (4) karena hanya memodifikasi sebagian kecil daftar rekomendasi (sesuai parameter $\alpha$) untuk menyuntikkan keberagaman, sementara sisa daftar tetap mempertahankan personalisasi berbasis kompetensi individu.
    
    \item \textbf{Kesimpulan:} Metode \textit{Constraint-based Re-ranking} dipilih. Metode \textit{re-ranking} ini menawarkan keseimbangan terbaik karena bekerja otomatis seperti FRediECH dalam memecahkan isolasi, namun fleksibel seperti UCRS dalam hal implementasi tanpa data (\textit{cold-start}). Sistem akan menggunakan algoritma ini untuk memfilter calon pasangan \textit{peer assessment}, guna memastikan setiap siswa terpapar pada rekan dengan tingkat kemahiran yang beragam (heterogen) sesuai batasan statistik yang ditentukan.
\end{enumerate}

\subsection{Keputusan Akhir Pemilihan Solusi}
\label{ssec:keputusan_akhir_pemilihan_solusi}

Berdasarkan analisis komparatif yang telah dilakukan melalui matriks keputusan, penelitian ini menetapkan arsitektur solusi yang terdiri dari tiga komponen utama yang saling terintegrasi. Pemilihan ini didasarkan pada prinsip efisiensi komputasi, ketersediaan data pada fase awal (\textit{cold-start}), dan validitas pedagogis untuk memitigasi risiko \textit{overpersonalization}.

Komponen pertama yang menjadi fondasi pelacakan kompetensi adalah \textbf{Sistem \textit{Elo rating}}. Berbeda dengan model probabilitas statis seperti IRT atau model berbasis \textit{Deep Learning} seperti DKT yang membutuhkan dataset historis masif, \textit{Elo rating} dipilih karena kemampuannya untuk beradaptasi secara dinamis dengan ukuran sampel yang kecil. Algoritma ini akan berfungsi sebagai mesin utama untuk mengestimasi parameter kemahiran siswa secara \textit{real-time} setelah setiap interaksi, memberikan landasan data yang akurat namun ringan secara komputasi untuk proses adaptasi selanjutnya \autocite{Vesin2022}.

Sebagai mekanisme interaksi sosial, sistem akan mengimplementasikan \textbf{\textit{Peer Assessment}} (penilaian sejawat). Keputusan ini diambil karena aktivitas ini menawarkan keuntungan ganda: secara teknis, ia menghasilkan artefak data kuantitatif (skor penilaian) yang terstruktur dan mudah diolah oleh algoritma; secara pedagogis, ia memberikan manfaat metakognitif tertinggi bagi siswa yang berperan sebagai penilai (\textit{assessor}) \autocite{Lu2012}. Aktivitas ini menjadi wadah strategis bagi sistem untuk menyuntikkan intervensi sosial yang terukur, menggantikan interaksi diskusi bebas yang sulit dikuantifikasi.

Komponen krusial yang mengikat kedua elemen di atas untuk memecahkan isolasi pembelajaran adalah algoritma \textbf{\textit{Constraint-based Re-ranking}}. Diadaptasi dari kerangka kerja \textcite{DeBiasio2023}, teknik ini dipilih sebagai strategi mitigasi utama karena keunggulannya dalam menangani masalah \textit{cold-start} dibandingkan pendekatan berbasis neural network seperti FRediECH. Sistem akan menggunakan algoritma ini untuk memfilter daftar calon pasangan \textit{peer assessment} dengan menerapkan batasan (\textit{constraint}) statistik $\alpha$. Batasan ini memastikan bahwa setiap siswa terpapar pada proporsi minimum rekan yang memiliki tingkat kemahiran berbeda (heterogen), yang dihitung berdasarkan jarak skor Elo mereka. Dengan demikian, sistem menjamin terjadinya diversifikasi perspektif dan pemecahan \textit{echo chamber} kompetensi secara otomatis sejak hari pertama operasional.

Penerapan mekanisme ini secara fundamental mengubah definisi "konten pembelajaran" dalam sistem. Umpan balik (\textit{feedback}) yang dihasilkan dari aktivitas penilaian sejawat diperlakukan sebagai \textbf{artefak pembelajaran dinamis}. Berbeda dengan soal latihan yang direkomendasikan secara personal (homogen) oleh algoritma \textit{Elo rating}, umpan balik dari rekan sejawat yang dipilih secara heterogen merepresentasikan \textbf{eksposur terhadap konten yang tidak terpersonalisasi}. Hal ini memaksa siswa untuk menghadapi perspektif, logika penyelesaian masalah, atau kritik yang berada di luar "zona nyaman" yang terpersonalisasi oleh algoritma sehingga secara efektif meminimalkan risiko \textit{overpersonalization} tanpa merusak jalur belajar (\textit{learning path}) individu.